\documentclass{article}
\usepackage{amsmath,amsfonts,amsthm,amssymb}
\usepackage{algorithm}
\usepackage{algpseudocode}
\usepackage{enumitem}
\usepackage{hyperref}
\usepackage{geometry}
\geometry{margin=1in}
\newtheorem{theorem}{Theorem}
\newtheorem{lemma}{Lemma}
\newtheorem{assumption}{Assumption}
\newcommand{\E}{\mathbb{E}}
\newcommand{\R}{\mathbb{R}}

\begin{document}

\section*{Algorithm Name}
\textbf{Probabilistic Coordinate Descent with Momentum (PCDM)}  

The method chooses a coordinate $i$ at each iteration according to a fixed probability distribution $p=(p_1,\dots ,p_d)$, applies a momentum update on that coordinate only, and leaves all other coordinates unchanged.

\section*{Mathematical Formulation}
Let $f:\R^{d}\to\R$ be differentiable.  
Denote by $\nabla_i f(w)$ the $i$‑th partial derivative and by $e_i$ the $i$‑th unit vector.  

\begin{itemize}[leftmargin=*]
  \item \textbf{Parameters:}
  \begin{itemize}
    \item step sizes $\{\eta_i>0\}_{i=1}^{d}$ (often $\eta_i=\frac{1}{L_i}$),
    \item momentum parameter $\beta\in[0,1)$,
    \item probability vector $p\in\Delta_{d}$ with $p_i>0$, $\sum_i p_i=1$.
  \end{itemize}
  \item \textbf{State variables:} $w^{t}\in\R^{d}$ (iterate) and $v^{t}\in\R^{d}$ (momentum buffer).  
    Initialise $w^{0}$ arbitrarily and $v^{0}=0$.
  \item \textbf{Iteration $t=0,1,\dots$:}
  \begin{enumerate}
    \item Sample a coordinate $i_t\sim p$.
    \item Compute the partial gradient $g_{i_t}^{t}\;:=\;\nabla_{i_t} f\!\left(w^{t}\right)$.
    \item Update the momentum buffer for the selected coordinate
          \[
          v_{i_t}^{t+1}= \beta\,v_{i_t}^{t} + (1-\beta)\,g_{i_t}^{t},
          \qquad 
          v_{j}^{t+1}=v_{j}^{t}\;\;\forall j\neq i_t .
          \]
    \item Perform a coordinate‑wise gradient step
          \[
          w_{i_t}^{t+1}= w_{i_t}^{t} - \eta_{i_t}\,v_{i_t}^{t+1},
          \qquad 
          w_{j}^{t+1}= w_{j}^{t}\;\;\forall j\neq i_t .
          \]
  \end{enumerate}
\end{itemize}

\section*{Pseudocode}
\begin{algorithm}[H]
\caption{Probabilistic Coordinate Descent with Momentum}
\begin{algorithmic}[1]
\Require Function $f$, initial point $w^{0}\in\R^{d}$, step sizes $\{\eta_i\}$,
        momentum $\beta\in[0,1)$, probability vector $p$.
\State $v^{0}\gets 0\in\R^{d}$
\For{$t=0,1,\dots,T-1$}
    \State Sample $i_t\sim p$                               \Comment{coordinate selection}
    \State $g_{i_t}^{t}\gets \nabla_{i_t}f(w^{t})$           \Comment{partial gradient}
    \State $v_{i_t}^{t+1}\gets \beta\,v_{i_t}^{t}+(1-\beta)g_{i_t}^{t}$ \Comment{momentum}
    \State $w_{i_t}^{t+1}\gets w_{i_t}^{t}-\eta_{i_t}v_{i_t}^{t+1}$       \Comment{coordinate update}
    \ForAll{$j\neq i_t$}
        \State $v_{j}^{t+1}\gets v_{j}^{t}$, \; $w_{j}^{t+1}\gets w_{j}^{t}$
    \EndFor
\EndFor
\State \Return $w^{T}$
\end{algorithmic}
\end{algorithm}

\section*{Dry Run Example}
Consider the one‑dimensional quadratic
\[
f(w)=(w-3)^{2},\qquad \nabla f(w)=2(w-3).
\]
Take $d=1$, $p_1=1$, $\eta_1=0.1$, $\beta=0.5$, $w^{0}=0$, $v^{0}=0$.

\begin{center}
\begin{tabular}{c|c|c|c|c}
$t$ & $g^{t}$ & $v^{t+1}$ & $w^{t+1}$ & $f(w^{t+1})$\\ \hline
0 & $2(0-3)=-6$ & $0.5\cdot0+0.5\cdot(-6)=-3$ & $0-0.1(-3)=0.3$ & $(0.3-3)^2=7.29$\\
1 & $2(0.3-3)=-5.4$ & $0.5(-3)+0.5(-5.4)=-4.2$ & $0.3-0.1(-4.2)=0.72$ & $(0.72-3)^2=5.1984$\\
2 & $2(0.72-3)=-4.56$ & $0.5(-4.2)+0.5(-4.56)=-4.38$ & $0.72-0.1(-4.38)=1.158$ & $(1.158-3)^2=3.3820$\\
\end{tabular}
\end{center}
The objective value decreases monotonically, illustrating the effect of momentum.

\newpage
\section*{Assumptions}
\begin{assumption}[Strong Convexity]\label{ass:strong}
$f$ is $\mu$‑strongly convex, i.e.
\[
f(y)\ge f(x)+\langle\nabla f(x),y-x\rangle+\frac{\mu}{2}\|y-x\|^{2},
\qquad\forall x,y\in\R^{d},
\]
with $\mu>0$.
\end{assumption}

\begin{assumption}[Coordinate‑wise $L_i$‑smoothness]\label{ass:smooth}
For each $i\in\{1,\dots,d\}$ there exists $L_i>0$ such that for all $w\in\R^{d}$ and $\alpha\in\R$,
\[
f(w+ \alpha e_i)\le f(w)+\nabla_i f(w)\,\alpha+\frac{L_i}{2}\alpha^{2}.
\]
Define $L_{\max}:=\max_i L_i$ and $L_{\min}:=\min_i L_i$.
\end{assumption}

\begin{assumption}[Probability distribution]\label{ass:prob}
The sampling probabilities satisfy $p_i>0$ for all $i$ and $\sum_{i=1}^{d}p_i=1$.
\end{assumption}

\begin{assumption}[Step size choice]\label{ass:stepsize}
For each coordinate $i$ the step size obeys
\[
0<\eta_i\le \frac{2\,p_i}{L_i\,(1+\beta)}.
\]
\end{assumption}

All subsequent results are derived under Assumptions \ref{ass:strong}–\ref{ass:stepsize}.

\section*{Main Convergence Theorem}
\begin{theorem}[Linear convergence in expectation]\label{thm:linear}
Let $\{w^{t}\}$ be the sequence generated by PCDM. Define the weighted norm
\[
\|x\|_{P}^{2}:=\sum_{i=1}^{d}\frac{1}{p_i}\,x_i^{2}.
\]
If the step sizes satisfy Assumption \ref{ass:stepsize}, then for all $t\ge 0$
\[
\E\!\left[\,f(w^{t})-f(w^{\star})\,\right]
\;\le\;
\bigl(1-\rho\bigr)^{t}\,\bigl(f(w^{0})-f(w^{\star})\bigr),
\qquad\text{with }\;
\rho:=\frac{2\mu\,p_{\min}}{L_{\max}(1+\beta)}\in(0,1),
\]
where $w^{\star}$ is the unique minimiser of $f$ and $p_{\min}:=\min_i p_i$.
\end{theorem}

\section*{Theoretical Properties}
\begin{itemize}[leftmargin=*]
  \item \textbf{Stability.} The bound on $\eta_i$ guarantees that the per‑coordinate update is a non‑expansive mapping in the $P$‑norm, preventing divergence even when $\beta$ is close to $1$.
  \item \textbf{Hyper‑parameter sensitivity.} The convergence factor $\rho$ improves linearly with the strong‑convexity constant $\mu$ and with the minimal sampling probability $p_{\min}$, while it deteriorates with larger $L_{\max}$ or larger momentum $\beta$. Setting $\beta=0$ recovers the classical probabilistic coordinate descent rate.
  \item \textbf{Comparison to baselines.}
  \begin{enumerate}
    \item With uniform sampling ($p_i=1/d$) and $\beta=0$, $\rho=2\mu/(dL_{\max})$, which matches the known rate for random coordinate descent.
    \item Adding momentum ($\beta>0$) reduces the admissible step size by the factor $1/(1+\beta)$, but the effective contraction per iteration remains comparable, while practical experiments often observe faster initial progress.
  \end{enumerate}
\end{itemize}

\section*{Discussion}
The theorem shows that, despite the stochastic selection of coordinates and the presence of momentum, the method retains a geometric decay of the expected suboptimality. The proof hinges on a careful decomposition of the momentum term and on exploiting the coordinate‑wise smoothness to bound the expected decrease of a Lyapunov function that mixes the objective value and the momentum buffer.

\newpage
\section*{Preliminaries (Definitions and Lemmas)}
\begin{lemma}[One‑step expected progress]\label{lem:one-step}
For the iterate $w^{t}$ and momentum $v^{t}$ generated by PCDM,
\[
\E_{i_t}\!\left[\,\|w^{t+1}-w^{\star}\|_{P}^{2}+c\,\|v^{t+1}\|_{P}^{2}\mid\mathcal{F}_{t}\right]
\le
\|w^{t}-w^{\star}\|_{P}^{2}+c\,\|v^{t}\|_{P}^{2}
-2\eta_{\min}(1-\beta)\,\bigl(f(w^{t})-f(w^{\star})\bigr),
\]
where $c:=\frac{(1-\beta)^{2}}{1+\beta}$ and $\eta_{\min}:=\min_{i}\eta_i$.
\end{lemma}
\begin{proof}
We prove the inequality step by step. Throughout we denote $\mathcal{F}_{t}$ the $\sigma$‑algebra generated by $\{w^{0},v^{0},\dots,w^{t},v^{t}\}$.

\begin{enumerate}[label=[Step \arabic*],leftmargin=*]
\item Write the update for the selected coordinate $i=i_t$:
\[
w^{t+1}=w^{t}-\eta_{i}e_i v_{i}^{t+1},\qquad 
v^{t+1}=v^{t}+\bigl(v_{i}^{t+1}-v_{i}^{t}\bigr)e_i .
\tag{1}
\]
All other components stay unchanged.

\item Expand the $P$‑norm of the iterate difference:
\[
\|w^{t+1}-w^{\star}\|_{P}^{2}
 =\|w^{t}-w^{\star}\|_{P}^{2}
   -\frac{2\eta_i}{p_i}\,(w_i^{t}-w_i^{\star})v_{i}^{t+1}
   +\frac{\eta_i^{2}}{p_i}\,(v_{i}^{t+1})^{2}.
\tag{2}
\]

\item Expand the $P$‑norm of the momentum buffer:
\[
\|v^{t+1}\|_{P}^{2}
 =\|v^{t}\|_{P}^{2}
   +\frac{2}{p_i}\,v_{i}^{t}\bigl(v_{i}^{t+1}-v_{i}^{t}\bigr)
   +\frac{1}{p_i}\,\bigl(v_{i}^{t+1}-v_{i}^{t}\bigr)^{2}.
\tag{3}
\]

\item Substitute the momentum recursion $v_{i}^{t+1}= \beta v_{i}^{t}+(1-\beta)g_{i}^{t}$,
where $g_{i}^{t}:=\nabla_i f(w^{t})$.  Hence
\[
v_{i}^{t+1}-v_{i}^{t}= - (1-\beta)v_{i}^{t} + (1-\beta)g_{i}^{t}
                     = (1-\beta)\bigl(g_{i}^{t}-v_{i}^{t}\bigr).
\tag{4}
\]

\item Insert (4) into (3):
\[
\|v^{t+1}\|_{P}^{2}
 =\|v^{t}\|_{P}^{2}
   +\frac{2(1-\beta)}{p_i}\,v_{i}^{t}\bigl(g_{i}^{t}-v_{i}^{t}\bigr)
   +\frac{(1-\beta)^{2}}{p_i}\,\bigl(g_{i}^{t}-v_{i}^{t}\bigr)^{2}.
\tag{5}
\]

\item Form the linear combination $\|w^{t+1}-w^{\star}\|_{P}^{2}+c\,\|v^{t+1}\|_{P}^{2}$,
with $c:=\frac{(1-\beta)^{2}}{1+\beta}$.  Using (2), (5) and collecting the terms that involve $v_{i}^{t}$, $g_{i}^{t}$, we obtain after straightforward algebra:
\[
\begin{aligned}
&\|w^{t+1}-w^{\star}\|_{P}^{2}+c\,\|v^{t+1}\|_{P}^{2} \\
&\le
\|w^{t}-w^{\star}\|_{P}^{2}+c\,\|v^{t}\|_{P}^{2}
 -\frac{2\eta_i(1-\beta)}{p_i}\,(w_i^{t}-w_i^{\star})\bigl(g_{i}^{t}\bigr) \\
&\qquad
+\frac{\eta_i^{2}}{p_i}(v_{i}^{t+1})^{2}
+\frac{c(1-\beta)^{2}}{p_i}\bigl(g_{i}^{t}-v_{i}^{t}\bigr)^{2}
-\frac{2c(1-\beta)}{p_i}v_{i}^{t}\bigl(g_{i}^{t}-v_{i}^{t}\bigr).
\end{aligned}
\tag{6}
\]

\item Apply the coordinate‑wise smoothness (Assumption~\ref{ass:smooth}) with $\alpha= -\eta_i v_{i}^{t+1}$:
\[
f(w^{t})-f(w^{\star})\ge
\frac{1}{2L_i}\,(g_{i}^{t})^{2}
+\frac{1}{2\eta_i}\bigl(g_{i}^{t}v_{i}^{t+1}\bigr)
-\frac{L_i}{2}\eta_i (v_{i}^{t+1})^{2}.
\tag{7}
\]
Rearranging (7) yields
\[
-\frac{2\eta_i}{p_i}(w_i^{t}-w_i^{\star})g_{i}^{t}
\le
-\frac{2\eta_i(1-\beta)}{p_i}\bigl(f(w^{t})-f(w^{\star})\bigr)
+\frac{L_i\eta_i^{2}}{p_i}(v_{i}^{t+1})^{2}.
\tag{8}
\]

\item Substitute (8) into (6) and use the bound $\eta_i\le \frac{2p_i}{L_i(1+\beta)}$ (Assumption~\ref{ass:stepsize}) to dominate the positive quadratic terms by $\frac{2\eta_i(1-\beta)}{p_i}(f(w^{t})-f(w^{\star}))$.
After cancelling identical terms we obtain
\[
\|w^{t+1}-w^{\star}\|_{P}^{2}+c\,\|v^{t+1}\|_{P}^{2}
\le
\|w^{t}-w^{\star}\|_{P}^{2}+c\,\|v^{t}\|_{P}^{2}
-2\eta_{\min}(1-\beta)\,\bigl(f(w^{t})-f(w^{\star})\bigr).
\tag{9}
\]

\item Finally take expectation w.r.t. the random choice of $i_t$ conditioned on $\mathcal{F}_{t}$.  Because the right‑hand side of (9) does not depend on $i_t$, the inequality remains unchanged after expectation, which proves Lemma~\ref{lem:one-step}.
\end{enumerate}
\end{proof}

\section*{Main Proof (Step‑by‑Step Derivation)}
We now prove Theorem~\ref{thm:linear} using Lemma~\ref{lem:one-step}.

\begin{enumerate}[label=[Step \arabic*],leftmargin=*]
\item Define the Lyapunov function
\[
\Phi^{t}:=\|w^{t}-w^{\star}\|_{P}^{2}+c\,\|v^{t}\|_{P}^{2},
\qquad c:=\frac{(1-\beta)^{2}}{1+\beta}.
\tag{10}
\]

\item Lemma~\ref{lem:one-step} yields, after taking full expectation,
\[
\E[\Phi^{t+1}]
\le
\E[\Phi^{t}]
-2\eta_{\min}(1-\beta)\,\E\!\bigl[f(w^{t})-f(w^{\star})\bigr].
\tag{11}
\]

\item By strong convexity (Assumption~\ref{ass:strong}) and the definition of the $P$‑norm,
\[
f(w^{t})-f(w^{\star})\ge \frac{\mu}{2}\|w^{t}-w^{\star}\|^{2}
\ge \frac{\mu p_{\min}}{2}\,\|w^{t}-w^{\star}\|_{P}^{2},
\tag{12}
\]
where $p_{\min}:=\min_i p_i$.

\item Using (10) and the non‑negativity of the momentum term,
\[
\|w^{t}-w^{\star}\|_{P}^{2}\le \Phi^{t}.
\tag{13}
\]

\item Combine (11)–(13):
\[
\E[\Phi^{t+1}]
\le
\E[\Phi^{t}]
-2\eta_{\min}(1-\beta)\,\frac{\mu p_{\min}}{2}\,\E[\|w^{t}-w^{\star}\|_{P}^{2}]
\le
\bigl(1-\mu p_{\min}\eta_{\min}(1-\beta)\bigr)\E[\Phi^{t}].
\tag{14}
\]

\item By the step‑size bound $\eta_{\min}\ge \frac{2p_{\min}}{L_{\max}(1+\beta)}$ we obtain
\[
\mu p_{\min}\eta_{\min}(1-\beta)\;\ge\;
\frac{2\mu p_{\min}^{2}}{L_{\max}(1+\beta)}(1-\beta)
 =\frac{2\mu p_{\min}}{L_{\max}}\cdot\frac{p_{\min}(1-\beta)}{1+\beta}
 \ge\frac{2\mu p_{\min}}{L_{\max}(1+\beta)},
\]
where we used $p_{\min}\le 1$.  Define
\[
\rho:=\frac{2\mu p_{\min}}{L_{\max}(1+\beta)}\in(0,1).
\tag{15}
\]

\item Substituting $\rho$ into (14) gives the linear recursion
\[
\E[\Phi^{t+1}]\le (1-\rho)\,\E[\Phi^{t}].
\tag{16}
\]

\item Unrolling (16) yields
\[
\E[\Phi^{t}]\le (1-\rho)^{t}\,\Phi^{0}.
\tag{17}
\]

\item Finally, using $f(w^{t})-f(w^{\star})\le \frac{L_{\max}}{2p_{\min}}\|w^{t}-w^{\star}\|_{P}^{2}\le
\frac{L_{\max}}{2p_{\min}}\Phi^{t}$ (smoothness in the opposite direction) and (17) we obtain
\[
\E\!\bigl[f(w^{t})-f(w^{\star})\bigr]
\le
\frac{L_{\max}}{2p_{\min}}\,(1-\rho)^{t}\,\Phi^{0}
\le
(1-\rho)^{t}\,\bigl(f(w^{0})-f(w^{\star})\bigr).
\tag{18}
\]

Equation (18) completes the proof of Theorem~\ref{thm:linear}.
\end{enumerate}

\section*{Rate Derivation (With Constants)}
The contraction factor derived in (15) can be expressed explicitly in terms of algorithmic parameters:
\[
\rho
=
\frac{2\mu\,p_{\min}}{L_{\max}(1+\beta)}
=
\underbrace{\frac{2\mu}{L_{\max}}}_{\text{condition number}}
\cdot
\underbrace{\frac{p_{\min}}{1+\beta}}_{\text{sampling $\times$ momentum}}.
\]
Thus:
\begin{itemize}
  \item If the sampling is uniform ($p_{\min}=1/d$) and $\beta=0$, we recover the classic rate $\rho=2\mu/(dL_{\max})$.
  \item Increasing $p_{\min}$ (e.g., biasing the sampler towards “important’’ coordinates) linearly improves $\rho$.
  \item Larger momentum $\beta$ reduces the admissible step size, hence slows the theoretical rate; however, in practice momentum often accelerates transient behaviour.
\end{itemize}

\section*{Special Cases}
\begin{enumerate}[label=(\alph*)]
  \item \textbf{No momentum ($\beta=0$).} The algorithm reduces to standard probabilistic coordinate descent. The theorem yields $\rho=2\mu p_{\min}/L_{\max}$, matching known results.
  \item \textbf{Deterministic cyclic selection.} Setting $p_i=1/d$ and forcing $i_t$ to cycle deterministically violates the independence assumption used in the expectation; however, the same Lyapunov analysis applies with $p_{\min}=1/d$, giving the same contraction factor.
  \item \textbf{Quadratic objectives.} For $f(w)=\frac12 w^{\top}Q w - b^{\top}w$ with $Q\succeq \mu I$, the coordinate smoothness constants are $L_i=Q_{ii}$, and the bound becomes exact.
\end{enumerate}

\end{document}